% Options for packages loaded elsewhere
\PassOptionsToPackage{unicode}{hyperref}
\PassOptionsToPackage{hyphens}{url}
\documentclass[
]{article}
\usepackage{xcolor}
\usepackage{amsmath,amssymb}
\setcounter{secnumdepth}{-\maxdimen} % remove section numbering
\usepackage{iftex}
\ifPDFTeX
  \usepackage[T1]{fontenc}
  \usepackage[utf8]{inputenc}
  \usepackage{textcomp} % provide euro and other symbols
\else % if luatex or xetex
  \usepackage{unicode-math} % this also loads fontspec
  \defaultfontfeatures{Scale=MatchLowercase}
  \defaultfontfeatures[\rmfamily]{Ligatures=TeX,Scale=1}
\fi
\usepackage{lmodern}
\ifPDFTeX\else
  % xetex/luatex font selection
\fi
% Use upquote if available, for straight quotes in verbatim environments
\IfFileExists{upquote.sty}{\usepackage{upquote}}{}
\IfFileExists{microtype.sty}{% use microtype if available
  \usepackage[]{microtype}
  \UseMicrotypeSet[protrusion]{basicmath} % disable protrusion for tt fonts
}{}
\makeatletter
\@ifundefined{KOMAClassName}{% if non-KOMA class
  \IfFileExists{parskip.sty}{%
    \usepackage{parskip}
  }{% else
    \setlength{\parindent}{0pt}
    \setlength{\parskip}{6pt plus 2pt minus 1pt}}
}{% if KOMA class
  \KOMAoptions{parskip=half}}
\makeatother
\usepackage{longtable,booktabs,array}
\newcounter{none} % for unnumbered tables
\usepackage{calc} % for calculating minipage widths
% Correct order of tables after \paragraph or \subparagraph
\usepackage{etoolbox}
\makeatletter
\patchcmd\longtable{\par}{\if@noskipsec\mbox{}\fi\par}{}{}
\makeatother
% Allow footnotes in longtable head/foot
\IfFileExists{footnotehyper.sty}{\usepackage{footnotehyper}}{\usepackage{footnote}}
\makesavenoteenv{longtable}
\setlength{\emergencystretch}{3em} % prevent overfull lines
\providecommand{\tightlist}{%
  \setlength{\itemsep}{0pt}\setlength{\parskip}{0pt}}
\usepackage{bookmark}
\IfFileExists{xurl.sty}{\usepackage{xurl}}{} % add URL line breaks if available
\urlstyle{same}
\hypersetup{
  hidelinks,
  pdfcreator={LaTeX via pandoc}}

\author{}
\date{}

\begin{document}

\section{The Event-Density Foundations of Quantum Field
Theory}\label{the-event-density-foundations-of-quantum-field-theory}

\subsection{A Unified Substrate-to-Continuum
Overview}\label{a-unified-substrate-to-continuum-overview}

\textbf{Allen Proxmire} \textbf{Collaborator:} Claude (AI collaborator)
\textbf{April 2026} \textbf{Series:} Event-Density Foundational Theorems
--- Program Overview

\begin{center}\rule{0.5\linewidth}{0.5pt}\end{center}

\subsection{Abstract}\label{abstract}

The Event Density (ED) program presents a substrate ontology in which
quantum mechanics, classical fluid dynamics, classical electromagnetism,
non-Abelian gauge theory, and the leading-order content of classical
gravitation arise as continuum-scale consequences of a single set of
substrate primitives. This overview integrates the closed
structural-foundation work of the program --- twenty-one forced theorems
plus the ED Combination Rule --- into a unified picture organized around
the substrate-to-continuum machinery established by the Diffusion
Coarse-Graining Theorem (DCGT). The unifying thread across QM, NS / MHD,
and Yang-Mills is the same: substrate primitives (chains, V1 / V5
participation kernels, T17 gauge rule-type, generalized minimal
coupling) coarse-grain under hydrodynamic-window scale separation
\(\ell_P \ll R_\mathrm{cg} \ll L_\mathrm{flow}\) to produce canonical-ED
dynamical content with form-FORCED structure and value-INHERITED
numerical content. The program supplies architectural classifications
under the ED-I-06 \emph{Fields and Forces} ontology (forces from
participation structures vs.~frame-kinematic frame artifacts
vs.~continuum-imposed constraints), Clay-relevance verdicts
(NS-Smoothness's Intermediate Path C for Clay-NS; the parallel YM-6
verdict for Clay-YM), and identification of load-bearing structural
conditions where each Clay-relevance verdict's positive content depends
on value-INHERITED quantities. ED does not solve the Clay Millennium
Problems, predict numerical coupling constants, derive the Standard
Model gauge group, or unify gravity with field theory at the
constructive-rigorous level. What ED \emph{does} deliver --- across all
sectors covered by the closed structural-foundation work --- is a
coherent substrate-grounded architectural account of the canonical-ED
dynamical content with explicit identification of the load-bearing
conditions for each sector's Clay-relevance verdict and an honest
demarcation between what is FORCED at substrate level and what is
INHERITED at value layer. This is the first program-level synthesis of
the ED foundational work; it is intended as a foundation for future arcs
(ED-10 spacetime emergence, Hall-MHD extension, substrate-gravity
extensions) rather than as a closed-arc completion.

\begin{center}\rule{0.5\linewidth}{0.5pt}\end{center}

\subsection{1. Introduction}\label{introduction}

\subsubsection{1.1 Motivation}\label{motivation}

Standard physics presents quantum mechanics, classical fluid dynamics,
classical electromagnetism, and non-Abelian gauge theory as separate
phenomenological frameworks, each with its own postulates, axioms, and
structural commitments. The Schrödinger equation is postulated; the
Navier-Stokes equation is phenomenological; Maxwell's equations are
derived from gauge invariance plus minimal coupling; Yang-Mills extends
Maxwell to non-Abelian gauge groups via the same mechanism but with
structural complications at the level of constructive existence. There
is no single underlying framework from which all four emerge.

The Event Density (ED) program presents such a framework. ED begins with
thirteen substrate primitives (chains, events, participation channels,
kernels, commitment structures, individuation, Lorentz covariance, and
related substrate-level structural commitments) and derives --- as
forced consequences --- the form of standard quantum mechanics (T1--T16,
Phase-1 closure 2026-04-26), the form of classical fluid dynamics (NS-2,
NS-Smoothness, NS-Turbulence, P4-NN, NS-Q), the form of classical
electromagnetism (T17, NS-MHD-1, NS-MHD-2), the form of non-Abelian
Yang-Mills theory (Arc YM, 2026-04-30), and the leading-order content of
classical gravitation (T19, T20, ED Combination Rule, T21, Theorem GR1).

This overview integrates these results. The unifying thread is the
substrate-to-continuum machinery established by the Diffusion
Coarse-Graining Theorem (DCGT, Arc D, closed 2026-04-30): under
hydrodynamic-window coarse-graining over substrate primitives, the
continuum-scale evolution equations emerge with form-FORCED structure
and value-INHERITED numerical content. The same machinery applies to
scalar diffusion, vector-field viscosity, gauge-field dynamics, and
minimal-coupling matter sources.

\subsubsection{1.2 Relationship to the NS Synthesis
Paper}\label{relationship-to-the-ns-synthesis-paper}

The Navier-Stokes Synthesis Paper (April 2026) presents the
architectural classification of NS / MHD content in the ED program with
five-part architecture: main text §3--§6 (NS-1 dimensional forcing, NS-2
form derivation, NS-Smoothness Intermediate Path C, NS-Turbulence
cascade analysis); Appendix C (MHD architectural classification +
ED-I-06 ontological reading); Appendix D (DCGT substrate-to-continuum
integration); Appendix E (Yang-Mills synthesis + Clay-relevance
statement).

The present overview \emph{includes} the NS Synthesis Paper's content as
its NS / MHD / YM sector and \emph{extends} it with the program's other
closed work --- Phase-1 QM emergence, T17--T18 / N1 / GR1 foundational
theorems, substrate gravity (T19 / T20 / ECR / T21), and the ED-I-06
ontological framework that runs across all sectors. The NS Synthesis
Paper is the canonical NS / MHD / YM publication; the present overview
is the canonical program-level publication.

\subsubsection{1.3 Relationship to the YM
Arc}\label{relationship-to-the-ym-arc}

The Yang-Mills arc (Arc YM-1 through YM-6, closed 2026-04-30) extends
the substrate-to-continuum machinery from the Abelian gauge sector
(Lorentz force in NS-MHD-2, derived in Arc D) to the non-Abelian sector
(continuum YM equation \(D_\mu F^{\mu\nu} = J^\nu\), mass-gap mechanism
from V1 finite width, OS-positivity audit at structural-suggestive
level, Clay-relevance statement parallel to NS-Smoothness's Intermediate
Path C). The YM arc is included in this overview as one of the
closed-arc applications of the substrate-to-continuum machinery,
alongside the NS / MHD work and the substrate-gravity preview.

\subsubsection{1.4 Relationship to the ED-I Ontological
Papers}\label{relationship-to-the-ed-i-ontological-papers}

ED-I-06 (\emph{Fields and Forces in Event Density}, Feb 2026) supplies
the ontological framework under which the canonical / non-ED structural
classifications across NS / MHD / YM uniformly read as instances of a
single boundary: forces sourced by stable participation structures
vs.~frame-kinematic frame artifacts vs.~continuum-imposed constraints.
ED-I-06 is referenced throughout this overview as the ontological roof;
it is not a forced theorem but functions as the
structural-classification template.

Other ED-I-* papers (e.g., ED-I-19 on the Fermi-polaron mass-gap
interpretation) are independent ED-program assets that interact
tangentially with the closed structural work; they are not load-bearing
for the present overview.

\begin{center}\rule{0.5\linewidth}{0.5pt}\end{center}

\subsection{2. Substrate Foundations}\label{substrate-foundations}

The closed structural-foundation theorems of the ED program, organized
by closure date:

\subsubsection{2.1 Phase-1 (QM emergence):
T1--T16}\label{phase-1-qm-emergence-t1t16}

Closed 2026-04-26 via twelve memos in
\texttt{arcs/quantum/foundations/}. Establishes that the four QM
postulates --- Hilbert-space structure, unitarity, observables,
measurement --- are FORCED-unconditional consequences of substrate
primitives. Specific load-bearing theorems include T1--T7
(substrate-to-Hilbert-space mapping), T8 (commutation relations),
T10--T16 (operator algebra, Schrödinger evolution, measurement rule).
Phase-1 closure is the QM-emergence arc's structural completion.

\subsubsection{2.2 Arc Q closure: T17
(Gauge-Fields-as-Rule-Type)}\label{arc-q-closure-t17-gauge-fields-as-rule-type}

Closed 2026-04-27 via nineteen memos in \texttt{arcs/arc-Q/}. Theorem 17
(FORCED-unconditional) states:

\begin{quote}
\emph{The gauge field \(A_\mu\) is the participation measure of a
structural rule-type \(\tau_g\) whose group / vertex / worldline /
vacuum content is forced across all four channels under a single unified
gauge-quotient identification; form FORCED by primitives
P-02/04/06/07/10/11/13 + Theorems 1--16; numerical magnitudes (specific
gauge group, coupling values, V1 kernel parameters, vacuum energy
density) INHERITED.}
\end{quote}

T17 supplies the substrate-level gauge structure (non-Abelian-capable
Killing-form-bearing Lie algebra), generalized minimal coupling vertex
(\(\partial_\mu \to \partial_\mu - igA_\mu^a T^a\)), and gauge-quotient
identification (clause C8). Required input for both NS-MHD (Lorentz
force) and Arc YM.

\subsubsection{2.3 Arc B closure: T18 (V1 Kernel
Retardation)}\label{arc-b-closure-t18-v1-kernel-retardation}

Closed 2026-04-27 via six memos in \texttt{arcs/arc-B/}. Theorem 18
(FORCED-unconditional) states:

\begin{quote}
\emph{The V1 vacuum response kernel is uniquely forced at primitive
level to have support restricted to the forward causal cone; no
symmetric, advanced, or hybrid kernel is constructible at primitive
level; the microscopic arrow of time is structurally FORCED at the
kernel level.}
\end{quote}

T18 supplies the forward-cone-only causal structure of the V1 kernel.
Required for OS-positivity preservation analysis (YM-5) and for the
analytic-structure preservation of substrate kernels under the continuum
limit.

\subsubsection{2.4 Arc N closure: Theorem N1 (V1 finite-width vacuum
kernel)}\label{arc-n-closure-theorem-n1-v1-finite-width-vacuum-kernel}

Closed 2026-04-24 via six memos in \texttt{arcs/arc-N/}. Theorem N1
(FORCED-unconditional) establishes V1's finite-width vacuum kernel
admissibility class --- bounded both ways by V1-δ (zero-width REFUTED)
and V1-∞ (infinite-width REFUTED) refutations. The V1 kernel has finite
spatial width \(\sim\ell_P\), smooth profile, and positive Fourier
transform. Required for R1 derivation (NS-Smooth-2 / Arc\_D\_4) and YM
substrate mass-gap mechanism (Arc\_YM\_3).

\subsubsection{2.5 Phase-3 closure: Theorem GR1 (V1 with Synge world
function)}\label{phase-3-closure-theorem-gr1-v1-with-synge-world-function}

Closed 2026-04-24 via six memos in \texttt{arcs/quantum/foundations/}.
Theorem GR1 (FORCED-unconditional) lifts V1 to curved spacetime via
Hadamard parametrix construction. Establishes the \emph{kernel-level}
gravitational-sector content of ED while leaving Einstein equations
themselves SPECULATIVE (GR-4A status). Provides the gravitational
analogue of T18's forward-cone causality.

\subsubsection{2.6 Substrate-Gravity Arc: T19 / T20 / ECR /
T21}\label{substrate-gravity-arc-t19-t20-ecr-t21}

Closed 2026-04-27 / 28 in \texttt{arcs/arc-SG/}. Four substrate-gravity
results:

\begin{itemize}
\tightlist
\item
  \textbf{T19 (Newton's law from substrate):} \(G = c^3\ell_P^2/\hbar\)
  derived from substrate primitives.
\item
  \textbf{T20 (transition acceleration):} \(a_0 = c\,H_0/(2\pi)\)
  derived from substrate primitives.
\item
  \textbf{ED Combination Rule:} \(a = \sqrt{a_N\,a_0}\) as the
  substrate-level composition rule for the cross-term between
  Newton-class and transition-class accelerations.
\item
  \textbf{T21 (slope-4 BTFR):} \(v^4 = G\,M\,a_0\), the baryonic
  Tully-Fisher relation slope, derived from T19 + T20 + ECR.
\end{itemize}

These are leading-order classical-gravitation results at substrate
level; they are pre-Einstein-equation content (the Einstein equation
GR-4A remains SPECULATIVE per ED-Phys-10 acoustic-metric guardrails).

\subsubsection{2.7 Arc D closure: DCGT (Diffusion Coarse-Graining
Theorem)}\label{arc-d-closure-dcgt-diffusion-coarse-graining-theorem}

Closed 2026-04-30 via six memos in \texttt{theory/Arc\_D/}. DCGT
(FORCED-unconditional) establishes the substrate-to-continuum bridge for
ED-QFT:

\begin{quote}
\emph{Under hydrodynamic-window coarse-graining over ED substrate
primitives (chains, V1 / V5 kernels, charged structural rule-types under
T17 minimal coupling), the continuum-scale evolution equations are:
(scalar) \(\partial_t\rho = \nabla\cdot(M(\rho)\nabla\rho)\);
(directional-field)
\(\rho\,\partial_t\mathbf{v} = \mu\nabla^2\mathbf{v} - \kappa\mu_\mathrm{V1}\ell_P^2\nabla^4\mathbf{v} + \lambda\partial_t\mathbf{v} + \rho_q\mathbf{E} + \mathbf{j}\times\mathbf{B} - \rho(\mathbf{v}\cdot\nabla)\mathbf{v} - \nabla p\).
Form FORCED, signs FORCED stabilizing by V1 positive Fourier transform,
values INHERITED.}
\end{quote}

DCGT joins T17 / T18 / T19 in the structural-foundation theorem
inventory as the substrate-to-continuum bridge theorem of the program.

\begin{center}\rule{0.5\linewidth}{0.5pt}\end{center}

\subsection{3. Substrate → Continuum Machinery
(DCGT)}\label{substrate-continuum-machinery-dcgt}

\subsubsection{3.1 Hydrodynamic window}\label{hydrodynamic-window}

A coarse-graining cell \(\Omega(\mathbf{x},R_\mathrm{cg})\) of radius
\(R_\mathrm{cg}\) is in the \emph{hydrodynamic window} iff
\(\ell_P \ll R_\mathrm{cg} \ll L_\mathrm{flow}\). Lower bound suppresses
substrate-discreteness fluctuations; upper bound preserves
macroscopic-flow gradients. Temporal analogue:
\(\tau_\mathrm{V5} \ll \tau_\mathrm{cg} \ll \tau_\mathrm{flow}\).

\subsubsection{3.2 Multi-scale expansion}\label{multi-scale-expansion}

The coarse-graining operator
\(\langle f\rangle_{R_\mathrm{cg}}(\mathbf{x}) = |\Omega|^{-1}\int_\Omega f(\mathbf{x}+\boldsymbol{\delta})\,d^3\boldsymbol{\delta}\)
is applied to substrate flux statistics (event flux \(\mathbf{J}_\rho\),
momentum-flux \(\Pi_{ij}\), charged-chain momentum-flux
\(\Pi^{(q)}_{ij}\), V5 temporal kernel). Even moments of V1 produce
successive Laplacian-class corrections; odd moments vanish by V1
isotropy; V5 temporal moments produce Maxwell-class memory.

\subsubsection{3.3 Scalar diffusion}\label{scalar-diffusion}

Cell-averaged participation density under V1-mediated chain-step
transition statistics →
\(\partial_t\rho = \nabla\cdot(M(\rho)\nabla\rho) + \mathcal{O}(\ell_P^2\nabla^4\rho)\).
Substrate-derived mobility \(M(\rho)\) has P4-class saturation
(\(M(\rho_\mathrm{max}) = 0\), \(M > 0\) in bulk, smooth, monotonically
decreasing toward packing). Form FORCED; values INHERITED.

\subsubsection{3.4 Directional-field
diffusion}\label{directional-field-diffusion}

Chain momentum-flux under V1 isotropy + first-order gradient expansion →
\(\rho\,\partial_t\mathbf{v} = \mu(\rho)\nabla^2\mathbf{v} - \rho(\mathbf{v}\cdot\nabla)\mathbf{v} - \nabla p + \mathcal{O}(\ell_P^2\nabla^4\mathbf{v})\).
Substrate-derived viscosity
\(\mu(\rho) = \tfrac{1}{3}\rho\bigl\langle\delta^2\bigr\rangle_{V1}\Gamma_0(\rho)\);
advection emerges as Eulerian bookkeeping of convective momentum flux
(frame-kinematic non-ED status preserved). Form FORCED; values
INHERITED.

\subsubsection{3.5 V1 → R1 hyperviscosity}\label{v1-r1-hyperviscosity}

V1 fourth-moment / second-moment ratio in the multi-scale expansion
produces \(-\kappa\mu_\mathrm{V1}\ell_P^2\nabla^4\mathbf{v}\). Sign
positive (stabilizing) by V1 positive-Fourier-transform property.
Matches NS-3.01 form-FORCED top-down derivation exactly. Same mechanism
in YM gives \(\ell_P^2\nabla^2 A\) (one derivative-order lower because
gauge-field kinetic term is already first-derivative).

\subsubsection{3.6 V5 → viscoelastic
memory}\label{v5-viscoelastic-memory}

V5 first temporal moment → Maxwell-class memory
\(\tau_R\dot\sigma + \sigma = 2\mu S\) with
\(\tau_R = \bigl\langle\tau\bigr\rangle_{V5}\). Substrate origin of
P4-NN-D5 viscoelastic ansatz. Form FORCED; values INHERITED.

\subsubsection{3.7 Minimal-coupling
coarse-graining}\label{minimal-coupling-coarse-graining}

T17 minimal coupling on charged chains →
\(\delta\Pi^{(q)}_{ij} = j_jA_i\) at fluid scale. Momentum-flux
divergence + cross-product identity + charge conservation →
\(\rho\,\partial_t\mathbf{v} \supset \rho_q\mathbf{E} + \mathbf{j}\times\mathbf{B}\)
(Abelian Lorentz force). Generalizes to non-Abelian via
\(\partial_\mu \to \partial_\mu - igA_\mu^a T^a\) + non-Abelian charge
conservation \(D_\mu J^\mu = 0\).

\subsubsection{3.8 DCGT theorem}\label{dcgt-theorem}

The aggregate substrate-to-continuum result. Form FORCED; sign-FORCED
stabilizing; value-INHERITED at coupling, kernel widths, and gauge-group
choice. Error bounds
\(\mathcal{O}((\ell_P/L_\mathrm{flow})^4) + \mathcal{O}((\tau_\mathrm{V5}/\tau_\mathrm{flow})^2)\).
Closes NS-2 + NS-MHD-2 + (via generalization) Arc YM substrate-bridge in
a single theorem.

\begin{center}\rule{0.5\linewidth}{0.5pt}\end{center}

\subsection{4. Quantum Mechanics Emergence
(Phase-1)}\label{quantum-mechanics-emergence-phase-1}

\subsubsection{4.1 Four QM postulates
FORCED}\label{four-qm-postulates-forced}

Phase-1 closure 2026-04-26 establishes that the four foundational
postulates of standard quantum mechanics --- Hilbert-space state
structure, unitary evolution, Hermitian observables, measurement
(probabilistic outcomes weighted by squared overlap) --- are
FORCED-unconditional consequences of ED substrate primitives, not
separate axiomatic commitments. T1--T7 establish the
substrate-to-Hilbert-space mapping; T8 establishes the canonical
commutation relations; T10--T16 establish the operator algebra,
Schrödinger evolution, and Born-rule measurement structure.

\subsubsection{4.2 Participation-structure interpretation of the
wavefunction}\label{participation-structure-interpretation-of-the-wavefunction}

The wavefunction \(\Psi(x)\) is the participation measure of a chain on
its participation channel: \(\Psi(x) = \sqrt{b(x)}\,e^{i\theta(x)}\)
where \(b(x)\) is the bandwidth (probability density) and \(\theta(x)\)
is the phase. Under the ED-I-06 ontology, \(\Psi\) is a
directional-field participation structure (orientation-bearing);
\(|\Psi|^2\) is the corresponding scalar-field participation density.
This reading is consistent with both the Born-rule probability
interpretation (scalar-field-class density) and the phase-coherence
interpretation (directional-field orientation).

\subsubsection{4.3 Operator algebra from
ED-gradients}\label{operator-algebra-from-ed-gradients}

Hermitian observables emerge as ED-gradient-class operators acting on
the participation channel. The canonical commutation relations
\([X, P] = i\hbar\) are FORCED at primitive level by the
substrate-discreteness + finite-proper-time-interval primitives (P-01 +
P-13). The standard QM operator algebra (position, momentum, angular
momentum, etc.) inherits structurally from the substrate-discreteness
regime.

\subsubsection{4.4 No-collapse measurement
rule}\label{no-collapse-measurement-rule}

Measurement in ED is not a collapse of the wavefunction but a commitment
event at the substrate level --- a chain commits to a specific
participation channel via the P-11 commitment-irreversibility primitive.
The Born rule emerges as the ED-statistical weight of commitment to each
available channel. This avoids the standard quantum-mechanical
``measurement problem'' (no separate collapse postulate is required);
the standard probabilistic rule is FORCED at substrate level via the
structural mechanism of commitment.

\begin{center}\rule{0.5\linewidth}{0.5pt}\end{center}

\subsection{5. NS / MHD Program}\label{ns-mhd-program}

\subsubsection{5.1 Architectural
classification}\label{architectural-classification}

Eleven content items in incompressible MHD classified under ED-I-06: 6
canonical-ED (viscous diffusion, magnetic diffusion, Lorentz force, R1,
\(\partial_t\mathbf{B}\), \(\nabla\cdot\mathbf{B}=0\)); 2
fluid-mechanical-additions / continuum-imposed constraints (pressure,
\(\nabla\cdot\mathbf{v}=0\)); 3 transport-kinematic non-ED (advection,
induction kinematic, Ohm kinematic). Three-angle convergence on
advection-as-non-ED (NS-2.08 architectural / NS-Smoothness dynamical /
NS-Turbulence spectral); parallel three-angle convergence on
induction-kinematic-as-non-ED (NS-MHD-3).

\subsubsection{5.2 Mobility channel}\label{mobility-channel}

Substrate-derived mobility / viscosity from chain-step momentum-flux
statistics under V1 isotropy. Field-type-agnostic vector extension of
canonical PDE: same operator structure for scalar density and
directional fields (per Architectural Canon Vector Extension memo).

\subsubsection{5.3 Lorentz force}\label{lorentz-force}

T17 minimal coupling on charged chains → fluid-level
\(\rho_q\mathbf{E} + \mathbf{j}\times\mathbf{B}\). Kinematic
\(\mathbf{v}\times\mathbf{B}\) component derived from the
minimal-coupling form rather than separately committed
(kinematic-coupling-pattern refinement of NS-MHD-2 §6.2).

\subsubsection{5.4 R1 hyperviscous
stabilization}\label{r1-hyperviscous-stabilization}

Form-FORCED \(-\kappa\mu_\mathrm{V1}\ell_P^2\nabla^4\mathbf{v}\) from V1
finite-width vacuum kernel. Top-down derivation (NS-3.01) and bottom-up
substrate derivation (Arc\_D\_4) reconciled. NS-Smooth-2 establishes
strict monotone gradient-norm Lyapunov decay
\(-\kappa\mu_\mathrm{V1}\ell_P^2\|\nabla^3\mathbf{v}\|_2^2\) in ED-only
NS.

\subsubsection{5.5 V5 viscoelastic
memory}\label{v5-viscoelastic-memory-1}

V5 cross-chain temporal memory → Maxwell-class viscoelastic ansatz at
fluid scale. Substrate origin of P4-NN-D5 ansatz. Universal Mobility Law
β = 1.72 ± 0.37 across 10 systems within 1σ of canonical β = 2.

\subsubsection{5.6 NS-Smoothness Intermediate Path
C}\label{ns-smoothness-intermediate-path-c}

ED supplies a real Clay-NS-relevant regularizing mechanism (R1
form-FORCED stabilization). Quantitative competition against
destabilizing super-Burnett terms is INHERITED on both sides. Advective
vortex-stretching
\(\int\boldsymbol{\omega}\cdot S\boldsymbol{\omega}\,dV\) is the unique
indefinite-sign contribution to the gradient-norm Lyapunov in 3D NS ---
the unique transport-kinematic obstruction to ED-style smoothness.

\subsubsection{5.7 NS Synthesis Paper}\label{ns-synthesis-paper}

Five-part publication: main text §3--§6 + Appendix C (MHD + ED-I-06) +
Appendix D (DCGT) + Appendix E (YM synthesis). Currently published as
\texttt{.md} + \texttt{.tex} + \texttt{.pdf} in
\texttt{papers/Navier\ Stokes\_Synthesis\_Paper/}. Canonical NS / MHD /
YM closed-arc reference for the program.

\begin{center}\rule{0.5\linewidth}{0.5pt}\end{center}

\subsection{6. Yang-Mills Program}\label{yang-mills-program}

\subsubsection{6.1 Substrate → continuum YM derivation
(YM-2)}\label{substrate-continuum-ym-derivation-ym-2}

Continuum YM equation \(D_\mu F^{\mu\nu} = J^\nu\) derived from ED
substrate primitives via DCGT-style multi-scale expansion of non-Abelian
gauge-field correlators + T17 generalized minimal coupling on charged
chains + non-Abelian rule-type bracket structure on the gauge
generators. Field strength
\(F_{\mu\nu}^a = \partial_\mu A_\nu^a - \partial_\nu A_\mu^a + gf^{abc}A_\mu^b A_\nu^c\);
non-Abelian commutator term FORCED by T17 clause C2 (Lie-algebra
structure of \(\tau_g\)). Bianchi identity preserved as algebraic
structural consequence.

\subsubsection{6.2 Mass-gap mechanism
(YM-3)}\label{mass-gap-mechanism-ym-3}

V1 finite-width second-moment expansion produces effective mass scale
\(m_\mathrm{eff}^2 \sim c_{V1}\ell_P^{-2}\) for non-Abelian gauge
fields. Mass term gauge-invariant kinetic-class (not Proca-class).
Non-Abelian self-interaction (quartic \(g^2 A^4\) stabilization)
distinguishes the gap-bearing case from Abelian gapless case.
Continuum-limit survival conditional on kernel-profile rescaling
\(c_{V1}(\ell_P)\,\ell_P^{-2} \to m_\mathrm{phys}^2 > 0\).

\subsubsection{6.3 Architectural classification
(YM-4)}\label{architectural-classification-ym-4}

Six content channels classified: 4 canonical-ED (kinetic,
self-interaction, matter source, higher-derivative correction) / 2
non-ED (gauge-fixing, Christoffel artifacts). YM dynamics fully
canonical ED with no transport-kinematic obstruction class ---
structurally cleaner than NS or MHD.

\subsubsection{6.4 OS-positivity audit
(YM-5)}\label{os-positivity-audit-ym-5}

Channel-by-channel structural-suggestive positive: kinetic
(positive-definite for compact \(G\) via Killing form); self-interaction
(positive Euclidean quartic for compact \(G\)); matter source
(non-negative bilinear conditional on matter-sector OS positivity);
higher-derivative (\(\tfrac{1}{2}c_{V1}\ell_P^2\|\nabla A\|^2 \ge 0\)
via V1 positive Fourier transform). OS-positivity preservation locus:
compact gauge group + V1 positive Fourier transform + kernel-profile
rescaling + matter-sector OS positivity.

\subsubsection{6.5 Clay-relevance statement
(YM-6)}\label{clay-relevance-statement-ym-6}

Parallel in form and honesty to NS-Smoothness's Intermediate Path C. ED
substrate provides structurally suggestive path toward constructively
stable YM continuum limit; mass-gap mechanism FORCED at substrate level
with conditional survival; OS positivity preserved channel-by-channel
under canonical-ED content; gauge-fixing obstruction reframed via
substrate gauge-quotient identification (T17 C8). No constructive proof
claimed; structural-positive verdict identifying the load-bearing
conditions for a positive Clay-relevance verdict.

\begin{center}\rule{0.5\linewidth}{0.5pt}\end{center}

\subsection{7. Canonical ED Content Across All
Fields}\label{canonical-ed-content-across-all-fields}

A unified table of canonical-ED content channels across the program's
closed work:

{\def\LTcaptype{none} % do not increment counter
\begin{longtable}[]{@{}
  >{\raggedright\arraybackslash}p{(\linewidth - 6\tabcolsep) * \real{0.2500}}
  >{\raggedright\arraybackslash}p{(\linewidth - 6\tabcolsep) * \real{0.2500}}
  >{\raggedright\arraybackslash}p{(\linewidth - 6\tabcolsep) * \real{0.2500}}
  >{\raggedright\arraybackslash}p{(\linewidth - 6\tabcolsep) * \real{0.2500}}@{}}
\toprule\noalign{}
\begin{minipage}[b]{\linewidth}\raggedright
Sector
\end{minipage} & \begin{minipage}[b]{\linewidth}\raggedright
Content Channel
\end{minipage} & \begin{minipage}[b]{\linewidth}\raggedright
Substrate Origin
\end{minipage} & \begin{minipage}[b]{\linewidth}\raggedright
Form
\end{minipage} \\
\midrule\noalign{}
\endhead
\bottomrule\noalign{}
\endlastfoot
QM & Wavefunction \(\Psi\) & Participation measure of chain on
participation channel & Directional + scalar \\
QM & Hermitian observables & ED-gradient-class operators & \(X\), \(P\),
etc. via P-01 + P-13 \\
QM & Schrödinger evolution & Unitary participation-channel evolution &
\(i\hbar\partial_t\Psi = H\Psi\) \\
NS & Viscous diffusion & V1 second moment + P4-modulated \(\Gamma_0\) &
\(\mu(\rho)\nabla^2\mathbf{v}\) \\
NS & R1 hyperviscosity & V1 fourth-moment / second-moment &
\(-\kappa\mu_\mathrm{V1}\ell_P^2\nabla^4\mathbf{v}\) \\
NS & Maxwell viscoelastic & V5 first temporal moment &
\(\tau_R\dot\sigma + \sigma = 2\mu S\) \\
Maxwell & Magnetic diffusion & Mobility channel on \(\mathbf{B}\) &
\(\eta\nabla^2\mathbf{B}\) \\
Maxwell & Field structure & T17 directly & \(\partial_t\mathbf{B}\),
\(\nabla\cdot\mathbf{B}=0\), etc. \\
Maxwell & Lorentz force & T17 minimal coupling on charged chains &
\(\rho_q\mathbf{E} + \mathbf{j}\times\mathbf{B}\) \\
YM & Kinetic term & V1 directional-field curvature &
\(\partial_\mu F^{\mu\nu}\) \\
YM & Self-interaction & T17 rule-type commutator &
\(gf^{abc}A_\mu^b F^{\mu\nu\,c}\) \\
YM & Matter source & T17 generalized minimal coupling &
\(J^{\nu\,a}\) \\
YM & Higher-derivative & V1 second-moment expansion &
\(c_{V1}\ell_P^2\nabla^2 A\) \\
Gravity & Newton's law & T19: \(G = c^3\ell_P^2/\hbar\) &
\(\mathbf{F} = -G\,Mm/r^2\) \\
Gravity & Transition acceleration & T20: \(a_0 = c\,H_0/(2\pi)\) &
\(a_0\) scale \\
Gravity & Composition rule & ED Combination Rule &
\(a = \sqrt{a_N\,a_0}\) \\
\end{longtable}
}

Every entry is form-FORCED by substrate primitives + closed
structural-foundation theorems; every numerical content (specific values
of \(\mu\), \(\eta\), \(g\), \(G\), \(a_0\), \(\rho_\mathrm{max}\),
etc.) is INHERITED at value layer. The unifying methodology ---
form-FORCED / value-INHERITED --- runs across every sector.

\begin{center}\rule{0.5\linewidth}{0.5pt}\end{center}

\subsection{8. Mass-Gap Mechanisms Across
ED-QFT}\label{mass-gap-mechanisms-across-ed-qft}

\subsubsection{8.1 NS R1 (non-Clay)}\label{ns-r1-non-clay}

R1 stabilization \(-\kappa\mu_\mathrm{V1}\ell_P^2\nabla^4\mathbf{v}\)
from V1 finite-width vacuum kernel produces strict monotone Lyapunov
decay in ED-only NS (without advection). In full 3D NS, the advective
vortex-stretching
\(\int\boldsymbol{\omega}\cdot S\boldsymbol{\omega}\,dV\) is the unique
indefinite-sign Lyapunov contribution --- the obstruction to ED-style
smoothness lies outside the canonical regularizing architecture. NS R1
is \emph{not} a mass gap in the QFT spectral sense; it is a
Lyapunov-stabilization mechanism for fluid-velocity gradient-norm decay.
Non-Clay-relevant for NS but Clay-NS-relevant in the Intermediate Path C
sense (ED supplies the regularizing mechanism + identifies the
obstruction).

\subsubsection{8.2 YM substrate mass-gap mechanism (conditional
survival)}\label{ym-substrate-mass-gap-mechanism-conditional-survival}

V1 finite-width second-moment expansion produces effective mass scale
\(m_\mathrm{eff}^2 \sim c_{V1}\ell_P^{-2}\) for non-Abelian gauge
fields. Mass term gauge-invariant kinetic-class. Non-Abelian quartic
stabilization preserves the gap against loop-correction remixing.
Continuum-limit survival conditional on kernel-profile rescaling
\(c_{V1}(\ell_P)\ell_P^{-2} \to m_\mathrm{phys}^2 > 0\). ED supplies the
mechanism FORCED at substrate level; physical mass-gap value INHERITED
via kernel-profile-rescaling.

\subsubsection{8.3 Structural parallels and
differences}\label{structural-parallels-and-differences}

\textbf{Parallels:} Both mechanisms originate in V1 finite-width vacuum
kernel substrate properties (V1 second moment for both; same Theorem N1
admissibility class; same sign-FORCED stabilizing property by V1
positive Fourier transform). Both produce stabilizing higher-derivative
corrections at fluid / continuum scale (\(\nabla^4\mathbf{v}\) in NS;
\(\nabla^2 A\) in YM, the latter one derivative-order lower because
gauge-field kinetic is already first-derivative). Both have
value-INHERITED status for their physical-scale quantities.

\textbf{Differences:} NS R1 stabilizes \emph{Lyapunov decay}
(fluid-velocity gradient-norm); YM mass-gap stabilizes \emph{spectral
gap} (gauge-field momentum-mode suppression). NS smoothness obstruction
is in the transport-kinematic sector (advection); YM does not have a
transport-kinematic obstruction class (per YM-4 classification). NS R1
supplies the \emph{regularizing mechanism}; YM substrate gap supplies
the \emph{gap mechanism itself}. The Clay-relevance verdicts
(Intermediate Path C) are parallel in form but address structurally
different Clay problems.

\begin{center}\rule{0.5\linewidth}{0.5pt}\end{center}

\subsection{9. Positivity Structures}\label{positivity-structures}

\subsubsection{9.1 QM positivity}\label{qm-positivity}

Hilbert-space inner product \(\langle\Psi|\Phi\rangle\) is
positive-definite by construction; the Hermitian-observable spectrum is
real; the Born-rule probability \(|\langle\Psi|\Phi\rangle|^2\) is
non-negative. QM positivity is FORCED at the level of the Phase-1-closed
structure (T1--T7); not a separate axiom.

\subsubsection{9.2 NS Lyapunov decay}\label{ns-lyapunov-decay}

Gradient-norm Lyapunov \(L = \tfrac{1}{2}\|\nabla\mathbf{v}\|_2^2\) in
ED-only NS satisfies
\(dL/dt = -\nu\|\nabla^2\mathbf{v}\|_2^2 - \kappa\mu_\mathrm{V1}\ell_P^2\|\nabla^3\mathbf{v}\|_2^2 \le 0\)
--- strict monotone decay (NS-Smooth-2). In full 3D NS, advective
vortex-stretching is the unique indefinite-sign contribution; ED
supplies the regularizing structure but not the closure of Clay-NS
smoothness.

\subsubsection{9.3 YM OS positivity}\label{ym-os-positivity}

Each canonical-ED content channel preserves OS reflection positivity at
structural-suggestive level: kinetic positive-definite for compact
\(G\); self-interaction quartic positive in Euclidean signature; matter
source non-negative bilinear conditional on matter-sector OS positivity;
higher-derivative non-negative via V1 positive Fourier transform.
Combined Euclidean action bounded below + reflection-positive at the
structural-suggestive level. Constructive rigor remains a separate
technical question (Gribov-class obstructions in standard gauge-fixing
analyses).

\subsubsection{9.4 Substrate-level conditions for
stability}\label{substrate-level-conditions-for-stability}

Across all sectors, the load-bearing substrate-level structural
conditions for positivity / stability are the same:

\begin{itemize}
\tightlist
\item
  \textbf{V1 positive Fourier transform} (FORCED at substrate level by
  Theorem N1 + Theorem 18). Required for R1 sign-FORCED stabilizing in
  NS, kinetic-term positivity in YM, higher-derivative positivity in YM.
\item
  \textbf{Compact gauge group} (for YM only --- INHERITED at value
  layer). Required for kinetic-term and quartic-term positivity in YM.
\item
  \textbf{Kernel-profile rescaling} (for YM only --- INHERITED at value
  layer). Required for mass-gap survival under continuum limit.
\item
  \textbf{Matter-sector OS positivity} (FORCED at substrate level via
  closed T1--T18 work). Required for matter-source OS positivity in YM
  and for QM positivity itself.
\end{itemize}

\begin{center}\rule{0.5\linewidth}{0.5pt}\end{center}

\subsection{10. Substrate Gravity
(Preview)}\label{substrate-gravity-preview}

Closed substrate-gravity content as of 2026-04-30:

\begin{itemize}
\tightlist
\item
  \textbf{T19 (Newton's law from substrate):} \(G = c^3\ell_P^2/\hbar\)
  derived from substrate primitives. Newton's law of universal
  gravitation is a leading-order classical-gravitation result FORCED at
  substrate level.
\item
  \textbf{T20 (transition acceleration):} \(a_0 = c\,H_0/(2\pi)\) ---
  the MOND-class transition scale derived from substrate primitives via
  the cosmological-horizon participation-density connection.
\item
  \textbf{ED Combination Rule:} \(a = \sqrt{a_N\,a_0}\) --- the
  substrate-level composition rule between Newton-class and
  transition-class accelerations. Geometric-mean composition reflecting
  the substrate participation-density-bias structure.
\item
  \textbf{T21 (slope-4 BTFR):} \(v^4 = G\,M\,a_0\) --- the baryonic
  Tully-Fisher relation slope-4 prediction, derived from T19 + T20 +
  ECR.
\item
  \textbf{Theorem GR1 (curvature-like participation structure):} V1
  lifted to curved spacetime via Hadamard parametrix (Phase-3 closure
  2026-04-24). Provides the gravitational-sector kernel structure but
  does not derive the Einstein equation (GR-4A SPECULATIVE per
  ED-Phys-10 acoustic-metric guardrails).
\end{itemize}

These results are \emph{pre-Einstein-equation}: they cover the
leading-order classical-gravitation content (Newton + a₀ + BTFR) but do
not derive general relativity at the field-equation level. The ED-10
spacetime-emergence arc would extend this to full curvature emergence
(substrate-to-continuum metric coarse-graining; pick up the
curvature-like-field thread from ED-I-06 §5; upgrade Newton + a₀ to
genuine spacetime curvature). ED-10 is a queued long-horizon arc;
speculative ratio is high; estimated 8--12 memos.

\begin{center}\rule{0.5\linewidth}{0.5pt}\end{center}

\subsection{11. Program State and Open
Directions}\label{program-state-and-open-directions}

\subsubsection{11.1 Closed arcs (post
2026-04-30)}\label{closed-arcs-post-2026-04-30}

\begin{itemize}
\tightlist
\item
  \textbf{Phase-1 QM emergence} (2026-04-26): T1--T16, four QM
  postulates FORCED.
\item
  \textbf{Arc Q (Gauge-Fields-as-Rule-Type)} (2026-04-27): T17
  FORCED-unconditional.
\item
  \textbf{Arc B (Kernel-Level Arrow of Time)} (2026-04-27): T18
  FORCED-unconditional.
\item
  \textbf{Arc N (V1 Finite-Width Vacuum Kernel)} (2026-04-24): Theorem
  N1 FORCED-unconditional.
\item
  \textbf{Phase-3 (Gravitational Sector)} (2026-04-24): Theorem GR1
  FORCED-unconditional; GR-4A SPECULATIVE.
\item
  \textbf{Substrate-Gravity Arc} (2026-04-27 / 28): T19, T20, ECR, T21
  derived from substrate.
\item
  \textbf{NS / MHD Program} (2026-04-30): 11+5 arcs closed; NS Synthesis
  Paper published with Appendices C / D / E.
\item
  \textbf{Arc D (DCGT)} (2026-04-30): Substrate-to-continuum bridge
  theorem FORCED-unconditional.
\item
  \textbf{Yang-Mills Arc} (2026-04-30): YM-1 through YM-6;
  structural-positive Clay-relevance verdict.
\end{itemize}

\subsubsection{11.2 Open arcs}\label{open-arcs}

\begin{itemize}
\tightlist
\item
  \textbf{ED-10 spacetime emergence} --- substrate-to-continuum metric
  coarse-graining; curvature-like-field thread from ED-I-06 §5; upgrade
  substrate-gravity Newton + a₀ to genuine spacetime curvature. Most
  ambitious; speculative ratio high. Estimated 8--12 memos.
\item
  \textbf{Hall-MHD extension} --- audit Hall term +
  electron-pressure-gradient term in generalized Ohm's law. Optional
  arc-extension scope. Estimated 3--4 memos.
\item
  \textbf{Substrate-gravity extensions} --- extending T19 / T20 / ECR /
  T21 beyond closed-arc baseline. Would interact with ED-10 if both
  pursued.
\item
  \textbf{ED-I interpretation papers} --- additional ED-I-* papers as
  appropriate (parallel to ED-I-06 ontological-roof, ED-I-19
  Fermi-polaron). Discretionary scope; user-led.
\end{itemize}

\subsubsection{11.3 Future publication
arcs}\label{future-publication-arcs}

\begin{itemize}
\tightlist
\item
  \textbf{ED-10 Spacetime Emergence Paper} (after ED-10 arc closure).
\item
  \textbf{Substrate-Gravity Foundations Paper} (extends current
  substrate-gravity content to curvature regime if ED-10 closes).
\item
  \textbf{ED Program Manifesto} (book-length consolidation of ED across
  QFT, gravity, and emergence; long-horizon).
\end{itemize}

\begin{center}\rule{0.5\linewidth}{0.5pt}\end{center}

\subsection{12. Conclusion}\label{conclusion}

The Event Density program has, as of 2026-04-30, achieved a unified
substrate-to-continuum architecture covering quantum mechanics (Phase-1
closure), classical fluid dynamics (NS / MHD program), classical
electromagnetism (T17, NS-MHD-1, NS-MHD-2), non-Abelian gauge theory
(Yang-Mills arc), and the leading-order content of classical gravitation
(substrate-gravity arc with T19 / T20 / ECR / T21).

\textbf{ED-QFT now has a unified substrate → continuum architecture.}
The same machinery --- DCGT (Arc D) + T17 generalized minimal coupling +
T18 forward-cone kernel retardation + ED-I-06 ontological framework ---
applies across every closed-arc sector. The form-FORCED /
value-INHERITED methodology runs uniformly across QM, NS / MHD, YM, and
substrate-gravity. The ED-I-06 forces-vs-frame-kinematic-vs-constraint
ontological boundary serves as the structural-classification template
for canonical / non-canonical content.

\textbf{Structural results across QM, NS / MHD, YM, and DCGT are
mutually consistent.} Every closed-arc result respects the program's
closed structural-foundation theorems; no contradictions or open
structural inconsistencies have surfaced; the Phase-1 QM-emergence
content is consistent with the substrate-to-continuum machinery used in
the applied arcs; the Maxwell / Lorentz / YM gauge-sector content is
consistent with the T17 substrate gauge structure; the substrate-gravity
preview content (T19 / T20 / ECR / T21) is consistent with the
ED-Phys-10 acoustic-metric guardrails of Phase-3 GR1.

\textbf{Clay-relevance statements exist for NS and YM} at the
Intermediate Path C structural-positive level. ED supplies real
structural mechanisms (R1 hyperviscous regularization for NS smoothness;
substrate-induced mass scale + non-Abelian quartic stabilization for YM
mass gap); identifies the load-bearing conditions for each
Clay-relevance verdict (advection-as-non-ED unique obstruction for NS;
kernel-profile rescaling for YM); does not claim Clay-problem solutions.
The honest framing is parallel across both sectors: structural mechanism
FORCED at substrate level; physical-scale verdict conditional on
value-INHERITED quantities.

\textbf{The ED-QFT overview is a foundation for future arcs.} ED-10
spacetime emergence would extend the curvature-like-field thread of
ED-I-06 §5 + Theorem GR1 into full spacetime-curvature emergence,
upgrading substrate-gravity Newton + a₀ to the field-equation level.
Hall-MHD extension would audit additional fluid-mechanical content
channels under DCGT. Substrate-gravity extensions would extend T19 / T20
/ ECR / T21 beyond the closed baseline. Each future arc draws on the
closed structural-foundation work catalogued here.

\textbf{Honest program-level disclaimers.} ED does not solve the Clay
Millennium Problems (NS smoothness or YM existence + mass gap). ED does
not predict numerical coupling constants (\(g\), \(G\), \(a_0\),
\(\rho_\mathrm{max}\), etc.) --- all specific numerical values are
INHERITED at value layer. ED does not derive the Standard Model gauge
group (\(SU(3)\times SU(2)\times U(1)\) is empirical per Arc Q.6). ED
does not unify gravity with field theory at the constructive-rigorous
level --- gravitational sector content is leading-order classical
(Newton + a₀ + BTFR) with the Einstein equation SPECULATIVE per
ED-Phys-10. ED does not produce a full quantum-gravitational theory; it
produces a substrate-grounded architectural account of the canonical-ED
dynamical content across QM, NS / MHD, YM, and the gravitational kernel
sector, with explicit identification of what is FORCED and what is
INHERITED at every level.

\textbf{What ED \emph{does} deliver.} A coherent substrate-grounded
architectural account of canonical-ED dynamical content across QM, NS /
MHD, YM, and the leading-order classical-gravitation sector.
Substrate-to-continuum machinery (DCGT) covering scalar fields, vector
fields, gauge fields, matter coupling, and substrate-cutoff
regularization. Architectural classifications under ED-I-06 with
consistent canonical / non-canonical boundary across all sectors.
Clay-relevance verdicts at the Intermediate Path C level for NS and YM.
Identification of load-bearing conditions for each closed-arc result.
Foundation for future arcs including ED-10 spacetime emergence.

This is the Event Density program's first program-level synthesis.
Forty-plus memos are integrated into a single coherent picture;
twenty-one forced theorems plus DCGT, ED Combination Rule, Theorem N1,
Theorem GR1 provide the structural backbone; the form-FORCED /
value-INHERITED methodology supplies the consistent epistemological
framing. The program is structurally complete for its closed-arc scope
and ready for the next phase of arc development.

\begin{center}\rule{0.5\linewidth}{0.5pt}\end{center}

\subsection{References}\label{references}

\begin{itemize}
\tightlist
\item
  \emph{Event Density: One Substrate, Three Domains --- A Structural
  Account of Quantum Mechanics, Galactic Gravity, and Universal
  Soft-Matter Mobility} (Proxmire 2026). Earlier program-overview.
\item
  \emph{The Architectural Foundations of Navier-Stokes in Event Density}
  (Proxmire 2026). NS / MHD / YM closed-arc reference paper with
  five-part architecture (main + Appendix C + D + E).
\item
  \emph{Universal Degenerate-Mobility Scaling in Crowded Soft Matter}
  (Proxmire 2026). Universal Mobility Law empirical-anchor paper.
\item
  \emph{ED Mobility Saturation Predicts Non-Newtonian Rheology}
  (Proxmire 2026). P4-NN companion paper.
\item
  \emph{The Primitive-Level Arrow of Time in Event Density: The Closure
  of Arc B and Theorem 18} (Proxmire 2026). T18 V1 retarded-kernel
  structural-origin paper.
\item
  \emph{Structural Foundations of ED-Substrate Gravity: Newton, the
  Transition Scale, the Combination Rule, and the Baryonic Tully-Fisher
  Relation} (Proxmire 2026). Substrate-gravity foundations paper
  covering T19 / T20 / ECR / T21.
\item
  \emph{A UV-Finite Quantum Field Theory from Event-Density Primitives}
  (Proxmire 2026). Arc Q.8 vacuum factorisation paper.
\item
  \emph{Gauge Fields as Forced Rule-Type Structure: Theorem 17}
  (Proxmire 2026). Arc Q closure paper.
\item
  \emph{Fields and Forces in Event Density} (Proxmire, Feb 2026,
  ED-I-06). Ontological framework paper.
\end{itemize}

\begin{center}\rule{0.5\linewidth}{0.5pt}\end{center}

\emph{The Event-Density Foundations of Quantum Field Theory: A Unified
Substrate-to-Continuum Overview. First program-level synthesis
integrating Phase-1 QM emergence (T1--T16), Arc Q (T17), Arc B (T18),
Arc N (Theorem N1), Phase-3 (Theorem GR1), Substrate-Gravity (T19 / T20
/ ECR / T21), Arc D (DCGT), NS / MHD program, and Yang-Mills arc into a
coherent architectural account of canonical-ED dynamical content across
quantum, fluid-mechanical, gauge, and gravitational-kernel sectors.
Form-FORCED / value-INHERITED methodology consistent across every
sector. Clay-relevance verdicts at Intermediate Path C level for NS and
YM. ED-I-06 ontological framework as classification template.
Substrate-to-continuum machinery established for scalar, vector, gauge,
and minimal-coupling matter content. Structural results mutually
consistent; no open contradictions; foundation in place for ED-10
spacetime-emergence and other future arcs. Honest program-level
disclaimers preserved: ED does not solve Clay problems, does not predict
coupling constants, does not derive Standard Model, does not unify
gravity at constructive-rigorous level. What ED does deliver:
substrate-grounded architectural account with explicit FORCED /
INHERITED demarcation across the closed-arc landscape of the program.}

\end{document}
